\begin{figure}[H]
\centering  
\[
\setlength{\arraycolsep}{2pt}
\begin{array}{ccccccccc}

\mathcolorbox{lightorange}{A[1]} 
&\oplus& \gradientmathbox{lightorange}{lightred}{A[3]} 
&\oplus& \gradientmathbox{lightorange}{lightred}{A[5]} 
&\oplus& \mathcolorbox{lightorange}{A[7]} 
&=& \mathcolorbox{lightorange}{100}\\[1em]

\gradientmathbox{lightorange}{lightred}{A[3]}
&\oplus& \gradientmathbox{lightorange}{lightred}{A[5]}
&\oplus& \gradientmathbox{lightred}{lightpink}{A[8]} 
&\oplus& \mathcolorbox{lightgray}{A[0]} 
&=& \mathcolorbox{lightred}{011}\\[1em]

\gradientmathbox{lightred}{lightpink}{A[8]}
&\oplus& \mathcolorbox{lightpink}{A[9]} 
&\oplus& \gradientmathbox{lightpink}{lightpurple}{A[13]} 
&\oplus& \mathcolorbox{lightpink}{A[14]}
&=& \mathcolorbox{lightpink}{110}\\[1em]

\mathcolorbox{lightpurple}{A[10]} 
&\oplus& \gradientmathbox{lightpink}{lightpurple}{A[13]} 
&\oplus& \mathcolorbox{lightpurple}{A[15]} 
&\oplus& \mathcolorbox{lightgray}{A[0]} 
&=& \mathcolorbox{lightpurple}{111}
\end{array}
\]
\caption{System of linear equations with variables in $\mathbb{F}^3_2$ induced by Figure~\ref{fig:csc592_bigdata_nubff}. Each $A[i]$ is the $i$-th entry in the 
backing array $A$ (zero-indexed). The keys are $\mathcolorbox{lightorange}{x}$, $\mathcolorbox{lightred}{y}$, $\mathcolorbox{lightpink}{w}$ and 
$\mathcolorbox{lightpurple}{z}$ with fingerprints $\mathcolorbox{lightorange}{100}$, $\mathcolorbox{lightred}{011}$, $\mathcolorbox{lightpink}{110}$ and 
$\mathcolorbox{lightpurple}{111}$, respectively. Array entries and fingerprints are marked with corresponding key color(s). The phantom entry is colored gray. 
There is one equation for each key. }
\label{fig:csc592_bigdata_nusoe}
\end{figure}